% GENERATED FILE. Edit canonical lessons, then run npm run generate:books.
% canonical-source-hash: fnv1a64:f019a0676e0a15c5
% canonical-lessons: ML-C487-nagarasabha, ML-C487-villeju-ophisu, ML-C487-talukku, ML-C487-kalaktar, ML-C487-salyam

\chapter{Local Government}
\label{ch:ml-local-government}
\hlchaptermodality{487}

\begin{chapteropening}
\textbf{By the end of this chapter:} \emph{I can say a municipality, a village office, a taluk, a district collector and a nuisance.}

Complete the last lesson of chapter 487: I can say a municipality, a village office, a taluk, a district collector and a nuisance.
\end{chapteropening}

\section[\texorpdfstring{\ml{നഗരസഭ} (nagarasabha)}{നഗരസഭ (nagarasabha)}]{\ml{നഗരസഭ} (nagarasabha) — a municipality}

\label{lesson:ML-C487-nagarasabha}



\begin{quote}
Before the new one: say the Malayalam for office hours, then the Malayalam for a working day.
\end{quote}

\subsection*{You'll want to know: \ml{നഗരസഭ}}
\textbf{\ml{നഗരസഭ}} — \emph{nagarasabha} — \textquotedblleft{}a municipality\textquotedblright{}.

\begin{quote}
\ml{നഗരസഭ} \ml{തെരുവുകളിൽ} \ml{പുതിയ} \ml{വിളക്കുകൾ} \ml{വച്ചു}.
\end{quote}

\begin{quote}
\emph{nagarasabha teruvukaḷil putiya viḷakkukaḷ vaccu.}
\end{quote}

\begin{quote}
The municipality put up new lamps in the streets.
\end{quote}

The council of a town; a village has a \ml{പഞ്ചായത്ത്}.

\begin{cousinweb}
\textbf{\ml{നഗരസഭ}} (\emph{nagarasabha}) is a Malayalam compound built from \ml{നഗരം}, a town, from Sanskrit \emph{nagara} (CDIAL 6924), and \ml{സഭ}, an assembly, from Sanskrit \emph{sabhā} (Monier-Williams).
\end{cousinweb}

\begin{grammarlens}[title={What you've built}]
One more word for getting around.
\end{grammarlens}

\subsection*{Your turn}
\begin{itemize}
\raggedright
  \item \emph{Say it:} \emph{nagarasabha}
  \item \emph{Say it:} \emph{nagarasabha}, once more
  \item \emph{Say it:} \emph{nagarasabha teruvukaḷil putiya viḷakkukaḷ vaccu.} — The municipality put up new lamps in the streets.
  \item \emph{Recall:} say the Malayalam for office hours, then the Malayalam for a working day, then say \emph{nagarasabha} again
\end{itemize}

\begin{tcolorbox}[breakable,colback=teal!4,colframe=teal!35!black,title={Before you move on}]
What does \ml{നഗരസഭ} mean? (\textquotedblleft{}A municipality\textquotedblright{}.) Say it once more.
\end{tcolorbox}
\section[\texorpdfstring{\ml{വില്ലേജ്} \ml{ഓഫീസ്} (villējŭ ōphīsŭ)}{വില്ലേജ് ഓഫീസ് (villējŭ ōphīsŭ)}]{\ml{വില്ലേജ്} \ml{ഓഫീസ്} (villējŭ ōphīsŭ) — a village office}

\label{lesson:ML-C487-villeju-ophisu}



\begin{quote}
Before the new one: say the Malayalam for a working day, then the Malayalam for a municipality.
\end{quote}

\subsection*{You'll want to know: \ml{വില്ലേജ്} \ml{ഓഫീസ്}}
\textbf{\ml{വില്ലേജ്} \ml{ഓഫീസ്}} — \emph{villējŭ ōphīsŭ} — \textquotedblleft{}a village office\textquotedblright{}.

\begin{quote}
\ml{സർട്ടിഫിക്കറ്റിനായി} \ml{ഞാൻ} \ml{വില്ലേജ്} \ml{ഓഫീസിൽ} \ml{പോയി}.
\end{quote}

\begin{quote}
\emph{sarṭṭiphikkaṟṟināyi ñān villējŭ ōphīsil pōyi.}
\end{quote}

\begin{quote}
I went to the village office for a certificate.
\end{quote}

The office of the smallest revenue area; land records and many certificates come from it.

\begin{cousinweb}
\textbf{\ml{വില്ലേജ്} \ml{ഓഫീസ്}} (\emph{villējŭ ōphīsŭ}) was borrowed from English \emph{village office}, the name of the local revenue office.
\end{cousinweb}

\begin{grammarlens}[title={What you've built}]
One more word for getting around.
\end{grammarlens}

\subsection*{Your turn}
\begin{itemize}
\raggedright
  \item \emph{Say it:} \emph{villējŭ ōphīsŭ}
  \item \emph{Say it:} \emph{villējŭ ōphīsŭ}, once more
  \item \emph{Say it:} \emph{sarṭṭiphikkaṟṟināyi ñān villējŭ ōphīsil pōyi.} — I went to the village office for a certificate.
  \item \emph{Recall:} say the Malayalam for a working day, then the Malayalam for a municipality, then say \emph{villējŭ ōphīsŭ} again
\end{itemize}

\begin{tcolorbox}[breakable,colback=teal!4,colframe=teal!35!black,title={Before you move on}]
What does \ml{വില്ലേജ്} \ml{ഓഫീസ്} mean? (\textquotedblleft{}A village office\textquotedblright{}.) Say it once more.
\end{tcolorbox}
\section[\texorpdfstring{\ml{താലൂക്ക്} (tālūkkŭ)}{താലൂക്ക് (tālūkkŭ)}]{\ml{താലൂക്ക്} (tālūkkŭ) — a taluk}

\label{lesson:ML-C487-talukku}



\begin{quote}
Before the new one: say the Malayalam for a municipality, then the Malayalam for a village office.
\end{quote}

\subsection*{You'll want to know: \ml{താലൂക്ക്}}
\textbf{\ml{താലൂക്ക്}} — \emph{tālūkkŭ} — \textquotedblleft{}a taluk\textquotedblright{}.

\begin{quote}
\ml{ഞങ്ങളുടെ} \ml{ജില്ലയിൽ} \ml{ആറ്} \ml{താലൂക്കുകൾ} \ml{ഉണ്ട്}.
\end{quote}

\begin{quote}
\emph{ñaṅṅaḷuṭe jillayil āṟŭ tālūkkukaḷ uṇṭŭ.}
\end{quote}

\begin{quote}
Our district has six taluks.
\end{quote}

Plural: \ml{താലൂക്കുകൾ}, \emph{tālūkkukaḷ}.

A district is split into taluks, and a taluk into villages.

\begin{cousinweb}
\textbf{\ml{താലൂക്ക്}} (\emph{tālūkkŭ}) is usually traced to Arabic \emph{taʿalluq}, \textquotedblleft{}a dependency\textquotedblright{}, which reached the south through Hindustani as a word for a revenue division.
\end{cousinweb}

\begin{grammarlens}[title={What you've built}]
One more word for getting around.
\end{grammarlens}

\subsection*{Your turn}
\begin{itemize}
\raggedright
  \item \emph{Say it:} \emph{tālūkkŭ}
  \item \emph{Say it:} \emph{tālūkkŭ}, once more
  \item \emph{Say it:} \emph{ñaṅṅaḷuṭe jillayil āṟŭ tālūkkukaḷ uṇṭŭ.} — Our district has six taluks.
  \item \emph{Recall:} say the Malayalam for a municipality, then the Malayalam for a village office, then say \emph{tālūkkŭ} again
\end{itemize}

\begin{tcolorbox}[breakable,colback=teal!4,colframe=teal!35!black,title={Before you move on}]
What does \ml{താലൂക്ക്} mean? (\textquotedblleft{}A taluk\textquotedblright{}.) Say it once more.
\end{tcolorbox}
\section[\texorpdfstring{\ml{കളക്ടർ} (kaḷakṭar)}{കളക്ടർ (kaḷakṭar)}]{\ml{കളക്ടർ} (kaḷakṭar) — a district collector}

\label{lesson:ML-C487-kalaktar}



\begin{quote}
Before the new one: say the Malayalam for a village office, then the Malayalam for a taluk.
\end{quote}

\subsection*{You'll want to know: \ml{കളക്ടർ}}
\textbf{\ml{കളക്ടർ}} — \emph{kaḷakṭar} — \textquotedblleft{}a district collector\textquotedblright{}.

\begin{quote}
\ml{വെള്ളപ്പൊക്കം} \ml{കാരണം} \ml{കളക്ടർ} \ml{പള്ളിക്കൂടങ്ങൾക്ക്} \ml{അവധി} \ml{നൽകി}.
\end{quote}

\begin{quote}
\emph{veḷḷappokkaṁ kāraṇaṁ kaḷakṭar paḷḷikkūṭaṅṅaḷkkŭ avadhi nalki.}
\end{quote}

\begin{quote}
Because of the flood, the collector gave the schools a holiday.
\end{quote}

The officer who runs a district; the title goes back to collecting the land revenue.

\begin{cousinweb}
\textbf{\ml{കളക്ടർ}} (\emph{kaḷakṭar}) was borrowed from English \emph{collector}, the officer in charge of a district.
\end{cousinweb}

\begin{grammarlens}[title={What you've built}]
One more word for getting around.
\end{grammarlens}

\subsection*{Your turn}
\begin{itemize}
\raggedright
  \item \emph{Say it:} \emph{kaḷakṭar}
  \item \emph{Say it:} \emph{kaḷakṭar}, once more
  \item \emph{Say it:} \emph{veḷḷappokkaṁ kāraṇaṁ kaḷakṭar paḷḷikkūṭaṅṅaḷkkŭ avadhi nalki.} — Because of the flood, the collector gave the schools a holiday.
  \item \emph{Recall:} say the Malayalam for a village office, then the Malayalam for a taluk, then say \emph{kaḷakṭar} again
\end{itemize}

\begin{tcolorbox}[breakable,colback=teal!4,colframe=teal!35!black,title={Before you move on}]
What does \ml{കളക്ടർ} mean? (\textquotedblleft{}A district collector\textquotedblright{}.) Say it once more.
\end{tcolorbox}
\section[\texorpdfstring{\ml{ശല്യം} (śalyaṁ)}{ശല്യം (śalyaṁ)}]{\ml{ശല്യം} (śalyaṁ) — a nuisance; a bother}

\label{lesson:ML-C487-salyam}



\begin{quote}
Before the new one: say the Malayalam for a taluk, then the Malayalam for a district collector.
\end{quote}

\subsection*{You'll want to know: \ml{ശല്യം}}
\textbf{\ml{ശല്യം}} — \emph{śalyaṁ} — \textquotedblleft{}a nuisance; a bother\textquotedblright{}.

\begin{quote}
\ml{രാത്രി} \ml{നായകളുടെ} \ml{ശല്യം} \ml{കാരണം} \ml{ഞങ്ങൾക്ക്} \ml{ഉറങ്ങാൻ} \ml{കഴിഞ്ഞില്ല}.
\end{quote}

\begin{quote}
\emph{rātri nāyakaḷuṭe śalyaṁ kāraṇaṁ ñaṅṅaḷkkŭ uṟaṅṅān kaḻiññilla.}
\end{quote}

\begin{quote}
Because of the dogs' nuisance at night, we could not sleep.
\end{quote}

What bothers you comes first: \ml{കൊതുകുശല്യം}, the trouble from mosquitoes.

\begin{cousinweb}
\textbf{\ml{ശല്യം}} (\emph{śalyaṁ}) is borrowed from Sanskrit \emph{śalya}, \textquotedblleft{}a dart, a thorn\textquotedblright{} (Monier-Williams): something that sticks in and will not leave you alone.
\end{cousinweb}

\begin{grammarlens}[title={What you've built}]
One more word for getting around.
\end{grammarlens}

\subsection*{Your turn}
\begin{itemize}
\raggedright
  \item \emph{Say it:} \emph{śalyaṁ}
  \item \emph{Say it:} \emph{śalyaṁ}, once more
  \item \emph{Say it:} \emph{rātri nāyakaḷuṭe śalyaṁ kāraṇaṁ ñaṅṅaḷkkŭ uṟaṅṅān kaḻiññilla.} — Because of the dogs' nuisance at night, we could not sleep.
  \item \emph{Recall:} say the Malayalam for a taluk, then the Malayalam for a district collector, then say \emph{śalyaṁ} again
\end{itemize}

\begin{tcolorbox}[breakable,colback=teal!4,colframe=teal!35!black,title={Before you move on}]
What does \ml{ശല്യം} mean? (\textquotedblleft{}A nuisance; a bother\textquotedblright{}.) Say it once more.
\end{tcolorbox}
